% =============================================================================
\chapter{Theory}
\label{ch:theory}
% =============================================================================

This chapter presents the element formulation, global assembly, solution
procedures, stress recovery, and the optimisation algorithm implemented in
ShepherdLab LW~FEA-OPT.

% -----------------------------------------------------------------------------
\section{Element Formulation}
\label{sec:element}
% -----------------------------------------------------------------------------

\subsection{Degrees of Freedom}

Each node carries three degrees of freedom (DOFs):
%
\begin{equation}
    \vect{u}_i = \begin{bmatrix} u_{x,i} \\ u_{y,i} \\ \theta_i \end{bmatrix},
\end{equation}
%
where $u_x$ and $u_y$ are translational displacements and $\theta$ is the
in-plane rotation.  A single element connecting nodes~$i$ and~$j$ therefore
has a $6\times1$ displacement vector:
%
\begin{equation}
    \vect{u}^{(e)} =
    \begin{bmatrix} u_{x,i} & u_{y,i} & \theta_i & u_{x,j} & u_{y,j} & \theta_j \end{bmatrix}\T.
\end{equation}

\subsection{Local Stiffness Matrix}

In the \emph{local} coordinate system aligned with the element axis, the
standard 2-D Euler--Bernoulli frame stiffness matrix is
%
\begin{equation}
\label{eq:k_local}
\mat{k}^{(e)} =
\begin{bmatrix}
 \frac{EA}{L} & 0 & 0 & -\frac{EA}{L} & 0 & 0 \\[4pt]
 0 & \frac{12EI}{L^3} & \frac{6EI}{L^2} & 0 & -\frac{12EI}{L^3} & \frac{6EI}{L^2} \\[4pt]
 0 & \frac{6EI}{L^2} & \frac{4EI}{L} & 0 & -\frac{6EI}{L^2} & \frac{2EI}{L} \\[4pt]
 -\frac{EA}{L} & 0 & 0 & \frac{EA}{L} & 0 & 0 \\[4pt]
 0 & -\frac{12EI}{L^3} & -\frac{6EI}{L^2} & 0 & \frac{12EI}{L^3} & -\frac{6EI}{L^2} \\[4pt]
 0 & \frac{6EI}{L^2} & \frac{2EI}{L} & 0 & -\frac{6EI}{L^2} & \frac{4EI}{L}
\end{bmatrix},
\end{equation}
%
where $E$ is the elastic modulus, $A = w\,t$ is the cross-sectional area,
$I = w\,t^3/12$ is the second moment of area for a rectangular section of
width~$w$ and thickness~$t$, and $L$ is the element length.

\subsection{Coordinate Transformation}

An element oriented at angle~$\alpha$ (measured counter-clockwise from the
global $x$-axis) is rotated into the global frame by
%
\begin{equation}
\label{eq:transform}
\mat{T}^{(e)} =
\begin{bmatrix}
\mat{R} & \mat{0} \\
\mat{0} & \mat{R}
\end{bmatrix}, \qquad
\mat{R} =
\begin{bmatrix}
\cos\alpha & \sin\alpha & 0 \\
-\sin\alpha & \cos\alpha & 0 \\
0 & 0 & 1
\end{bmatrix}.
\end{equation}
%
The global element stiffness matrix is then
%
\begin{equation}
\label{eq:k_global}
    \mat{K}^{(e)} = {\mat{T}^{(e)}}\T\, \mat{k}^{(e)}\, \mat{T}^{(e)}.
\end{equation}

\subsection{Rigid Elements}

Rigid elements (e.g.\ attachment fixtures) are modelled using a stiffness
penalty method: the elastic modulus is multiplied by a large factor
$\beta$ (default $10^2$), giving $E_{\text{rigid}} = \beta\, E$.  This
is simpler than Lagrange multiplier approaches and integrates directly
with the standard assembly procedure.

\subsection{Cross-Section Properties}

All elements use a rectangular cross-section:
%
\begin{equation}
    A = w\,t, \qquad I = \frac{w\,t^3}{12},
\end{equation}
%
where $w$ is the (constant) out-of-plane width and $t$ is the thickness.
During optimisation, only $t$ varies; $w$ remains fixed.

% -----------------------------------------------------------------------------
\section{Global Assembly and Solution}
\label{sec:assembly}
% -----------------------------------------------------------------------------

\subsection{Assembly}

The global stiffness matrix is assembled from element contributions by
the standard direct stiffness method:
%
\begin{equation}
    \mat{K} = \sum_{e=1}^{M} \mat{L}^{(e)\mathsf{T}}\, \mat{K}^{(e)}\, \mat{L}^{(e)},
\end{equation}
%
where $\mat{L}^{(e)}$ is the Boolean localisation (scatter) matrix that maps
the six element DOFs to the global DOF vector.  In practice this is
implemented as index-based accumulation rather than explicit matrix
multiplication.

\subsection{Partitioning}

The global system $\mat{K}\,\vect{u} = \vect{P}$ is partitioned into free
(subscript~$f$) and fixed (subscript~$c$) DOFs:
%
\begin{equation}
\begin{bmatrix}
\mat{K}_{ff} & \mat{K}_{fc} \\
\mat{K}_{cf} & \mat{K}_{cc}
\end{bmatrix}
\begin{bmatrix} \vect{u}_f \\ \vect{u}_c \end{bmatrix}
=
\begin{bmatrix} \vect{P}_f \\ \vect{R}_c \end{bmatrix}.
\end{equation}
%
For prescribed zero displacements ($\vect{u}_c = \vect{0}$), the reduced
system is
%
\begin{equation}
    \mat{K}_{ff}\, \vect{u}_f = \vect{P}_f,
\end{equation}
%
which is solved by direct factorisation (MATLAB backslash operator).

% -----------------------------------------------------------------------------
\section{Newton--Raphson Large-Deflection Analysis}
\label{sec:newton_raphson}
% -----------------------------------------------------------------------------

For problems with large displacements, the equilibrium equation becomes
nonlinear because the element geometry (lengths and orientations) changes
with deformation.  LW~FEA-OPT uses an incremental load-stepping scheme
with Newton--Raphson iterations at each step.

\subsection{Load Stepping}

The total load $\vect{P}$ is divided into $N$ equal increments:
%
\begin{equation}
    \vect{P}_n = \frac{n}{N}\, \vect{P}_{\text{ref}}, \qquad n = 1, 2, \ldots, N.
\end{equation}

\subsection{Equilibrium and Residual}

At load step~$n$, the goal is to find displacements $\vect{u}$ such that
the internal forces balance the external forces:
%
\begin{equation}
    \vect{R}(\vect{u}) = \vect{f}_{\text{ext}} - \vect{f}_{\text{int}}(\vect{u}) = \vect{0}.
\end{equation}
%
The internal force vector is assembled element-by-element:
%
\begin{equation}
    \vect{f}_{\text{int}} = \sum_{e=1}^{M}
    {\mat{T}^{(e)}}\T\, \mat{k}^{(e)}\, \mat{T}^{(e)}\, \vect{u}^{(e)},
\end{equation}
%
where $\mat{T}^{(e)}$ and $\mat{k}^{(e)}$ are evaluated at the
\emph{current} deformed configuration.

\subsection{Iteration}

At each Newton--Raphson iteration~$k$:
%
\begin{enumerate}[nosep]
    \item Update element geometry (lengths, orientations) from the current
          displacement field.
    \item Assemble the tangent stiffness matrix
          $\mat{K}_t = \sum_e {\mat{T}^{(e)}}\T \mat{k}^{(e)} \mat{T}^{(e)}$.
    \item Compute the residual $\vect{R}^{(k)} = \vect{f}_{\text{ext}} - \vect{f}_{\text{int}}^{(k)}$.
    \item Solve for the displacement increment:
          $\mat{K}_{t,ff}\, \Delta\vect{u}_f = \vect{R}_f^{(k)}$.
    \item Update: $\vect{u}_f^{(k+1)} = \vect{u}_f^{(k)} + \Delta\vect{u}_f$.
\end{enumerate}

\subsection{Convergence Criteria}

Three criteria are checked simultaneously:
%
\begin{align}
    \text{Force:} \quad & \frac{\|\vect{R}\|}{\|\vect{f}_{\text{ext}}\|} < \varepsilon_f, \\[4pt]
    \text{Displacement:} \quad & \frac{\|\Delta\vect{u}\|}{u_{\max}} < \varepsilon_u, \\[4pt]
    \text{Energy:} \quad & |\vect{R}\T \Delta\vect{u}| < \varepsilon_E.
\end{align}
%
Convergence is declared when the force criterion is satisfied \emph{and}
either the displacement or energy criterion is met.  On the first iteration,
only the force criterion is used (the displacement increment has no
meaningful predecessor).

The default tolerances are $\varepsilon_f = 10^{-2}$,
$\varepsilon_u = 10^{-6}$, and $\varepsilon_E = 10^{-6}$.

% -----------------------------------------------------------------------------
\section{Stress Recovery}
\label{sec:stress}
% -----------------------------------------------------------------------------

After solving for displacements, member end forces are recovered in local
coordinates:
%
\begin{equation}
    \vect{f}_{\text{local}}^{(e)} = \mat{k}^{(e)}\, \mat{T}^{(e)}\, \vect{u}^{(e)},
\end{equation}
%
giving force components
$\vect{f}_{\text{local}} = [N_i,\; V_i,\; M_i,\; N_j,\; V_j,\; M_j]\T$
at the start ($i$) and end ($j$) nodes.

Stresses at each node of an element are:
%
\begin{align}
    \sigma_{\text{axial}}   &= \frac{N}{A}, \\[3pt]
    \sigma_{\text{bending}} &= \frac{M \cdot c}{I}, \quad c = t/2, \\[3pt]
    \tau_{\text{shear}}     &= \frac{3V}{2A},
\end{align}
%
where the shear stress uses the rectangular-section maximum (at the neutral
axis).  The combined von~Mises stress is
%
\begin{equation}
    \sigma_{\text{vM}} = \sqrt{(\sigma_{\text{axial}} + \sigma_{\text{bending}})^2
    + 3\,\tau_{\text{shear}}^2\,}.
\end{equation}

The factor of safety (FOS) for element~$e$ is
%
\begin{equation}
    \text{FOS}^{(e)} = \frac{\sigma_{\text{ult}}}
    {\max\!\big(\sigma_{\text{vM},i}^{(e)},\; \sigma_{\text{vM},j}^{(e)}\big)}.
\end{equation}

% -----------------------------------------------------------------------------
\section{Stress-Ratio Thickness Optimisation}
\label{sec:optimisation}
% -----------------------------------------------------------------------------

The optimiser implements a \emph{fully stressed design} (FSD) approach, which
iteratively adjusts element thicknesses so that every element operates near
its allowable stress.  The result is a more uniform stress distribution and,
typically, a lighter structure.

\subsection{Update Rule}

At iteration~$k$, the thickness of element~$e$ is updated as
%
\begin{equation}
\label{eq:fsd_update}
    t^{(k+1)}_e = t^{(k)}_e \left( \frac{\sigma^{(k)}_e}{\sigma_{\text{allow}}} \right)^{\!\alpha},
\end{equation}
%
where $\sigma^{(k)}_e$ is the maximum von~Mises stress in element~$e$ at
the current iteration, $\sigma_{\text{allow}} = \sigma_{\text{ult}} / \text{FOS}_{\text{target}}$,
and $\alpha \in (0,\,1]$ is a damping exponent that controls aggressiveness.

\begin{notebox}[Choosing $\alpha$]
Smaller values of $\alpha$ produce more cautious updates and are less
likely to oscillate, at the cost of slower convergence.  A typical range
is $0.3$--$0.5$.  For stiff structures, $\alpha = 0.4$ works well; for
flexible trusses, $\alpha = 0.3$ may be more stable.
\end{notebox}

\subsection{Move Limits}

To prevent oscillation, the relative change per iteration is clamped:
%
\begin{equation}
    t^{(k+1)}_e \in \Big[\, t^{(k)}_e (1 - \delta),\;\; t^{(k)}_e (1 + \delta) \,\Big],
\end{equation}
%
where $\delta$ is the move limit (default~$0.2$, i.e.\ $\pm 20\%$ per
iteration).

\subsection{Absolute Bounds}

After applying the move limit, thickness is clamped to user-specified bounds:
%
\begin{equation}
    t^{(k+1)}_e \in [\, t_{\min},\; t_{\max} \,].
\end{equation}
%
Setting a small $t_{\min}$ (e.g.\ $0.2$~mm) allows elements to effectively
vanish, enabling topology-like optimisation where load paths emerge
organically.

\subsection{Convergence}

Convergence is declared when the maximum relative thickness change across
all optimisable elements falls below a tolerance~$\varepsilon_t$:
%
\begin{equation}
    \max_e \frac{|t^{(k+1)}_e - t^{(k)}_e|}{t^{(k)}_e} < \varepsilon_t.
\end{equation}

\subsection{Stress Uniformity}

The quality of the optimised design is measured by the coefficient of
variation (CoV) of the maximum element stresses:
%
\begin{equation}
    \text{CoV} = \frac{\text{std}(\sigma_e)}{\text{mean}(\sigma_e)}.
\end{equation}
%
A lower CoV indicates more uniform material utilisation.

% -----------------------------------------------------------------------------
\section{Truss Subdivision}
\label{sec:subdivision}
% -----------------------------------------------------------------------------

To enable topology-like optimisation of truss structures, each original
member can be subdivided into $n$ smaller frame elements using
\texttt{subdivideTruss}.  This creates $n-1$ intermediate nodes per member,
uniformly spaced along the original member axis.

The fine mesh then allows the thickness optimiser to vary material
\emph{along} each member, not just between members.  Combined with a small
$t_{\min}$, the optimiser can effectively remove material from low-stress
regions, producing organic, topology-optimised shapes.

The \texttt{member\_map} output tracks which original member each
sub-element belongs to, which is useful for grouping, colouring, and
interpreting results.